% Method obstruction only. Not promoted to the main paper or counted as A5 movement.
\subsection*{A proper definable product bundle with no finite ambient Thom stratification}

Let $R$ be any real closed field. Set $Y=R^2$, $I=[-1,1]$, and define
the semialgebraic function
\[
 g(a,b,z)=
 \begin{cases}
 \displaystyle\frac{b^3(bz-a)}{b^4+(bz-a)^2},&(a,b)\ne(0,0),\\
 0,&(a,b)=(0,0),
 \end{cases}
 \qquad (a,b,z)\in Y\times I.
 \tag{MT.1}
\]
Let $X$ be its graph in $R^4$, with coordinates $(a,b,z,w)$, and let
$f:X\to Y$ be the restriction of $(a,b,z,w)\mapsto(a,b)$.
We claim that $f$ is a globally semialgebraically trivial, definably
proper bundle with interval fibre, but has no finite definable $C^1$
Thom $a_f$ stratification in this given ambient embedding.

\paragraph{The exact product and properness.}
Put $c=bz-a$ and $D=b^4+c^2$. The denominator is positive away
from $(a,b)=(0,0)$. The identity
\[
 b^2D^2-4b^6c^2=b^2(b^4-c^2)^2\ge0
\]
gives $|g(a,b,z)|\le |b|/2$ wherever the rational expression is
defined. Thus the displayed extension is continuous at every point
$(0,0,z)$, uniformly in $z$. At all other points it is rational with
positive denominator, and hence continuous and $C^\infty$.
The map
\[
 H:Y\times I\longrightarrow X,\qquad
 H(a,b,z)=(a,b,z,g(a,b,z))
 \tag{MT.2}
\]
is a semialgebraic homeomorphism, with inverse the coordinate
projection onto $(a,b,z)$, and $fH=\operatorname{pr}_Y$.
For every closed bounded definable $K\subset Y$, the inverse
image $f^{-1}(K)$ is closed and bounded: it is the graph of the
continuous function on $K\times I$, and the preceding estimate
bounds its fourth coordinate. Hence $f$ is definably proper.
Over the real numbers this also proves ordinary properness.
Its maximal open manifold-and-local-product locus is all of $Y$.

\paragraph{Failure for every finite ambient Thom stratification.}
We use the standard convention: source and target are partitioned
into finitely many definable $C^1$ embedded manifolds, with the
frontier condition; each source stratum maps submersively to one
target stratum; and for every incident source pair $S,T$, limits
of the tangent kernels along $S$ contain the tangent kernel along
$T$. We allow arbitrary finite compatible refinements of both
source and target. In particular, the central fibre need not be
assumed in advance to be a union of strata.

Write
\[
 E=\{(0,0,z,0):-1<z<1\}.
\]
Finiteness and definable dimension imply that some source stratum
$T$ has intersection with $E$ containing a nonempty relatively
open interval. Choose $q=(0,0,z_0,0)$ in that interval. A small
segment of $E$ through $q$ lies in $T$, so
\[
 e=(0,0,1,0)\in\ker d(f|_T)_q.
 \tag{MT.3}
\]
Now approach $q$ by the semialgebraic curve
\[
 q_t=(tz_0,t,z_0,0),\qquad t>0.
 \tag{MT.4}
\]
It lies in $X$, since $bz-a=0$ there. Away from the central
fibre, $X$ is the smooth graph of $g$ near these points, and
$f$ is a submersion. Differentiation gives
\[
 \frac{\partial g}{\partial z}
 =\frac{b^4\bigl(b^4-(bz-a)^2\bigr)}
              {\bigl(b^4+(bz-a)^2\bigr)^2}.
\]
This derivative is exactly $1$ at every $q_t$. Consequently the
full vertical tangent line of $f$ at $q_t$ is the fixed line
\[
 L_\mathrm{diag}=R(0,0,1,1).
 \tag{MT.5}
\]
By o-minimality, after shortening the positive interval the curve
$q_t$ lies in one source stratum $S$. The kernel of
$d(f|_S)_{q_t}$ is a subspace of $L_\mathrm{diag}$. It has constant
dimension because $f|_S$ is a submersion to a single target stratum.
It is therefore either the zero space at every point of this tail,
or the fixed line $L_\mathrm{diag}$ at every point. Its limit $L$
exists and satisfies $L\subseteq L_\mathrm{diag}$.

The strata $S$ and $T$ cannot coincide. If they did, the constant
rank kernel bundle of the $C^1$ submersion on that stratum would be
continuous at $q$; its limit along $q_t$ would contain the nonzero
vector $e$ in (MT.3). But $e\notin L_\mathrm{diag}$.
Thus $S\ne T$ and $q\in T\cap\overline S$. The frontier condition
gives $T\subset\overline S$. Thom's condition would require
\[
 Re\subseteq\ker d(f|_T)_q\subseteq L
                 \subseteq L_\mathrm{diag},
\]
which is impossible. This proves the claim for every finite
definable stratification, not merely for an unrefined central fibre.

\paragraph{Closed manifold fibres and the remaining A.5 dimensions.}
The same obstruction has a version with closed $2$-manifold fibres
and $3$-dimensional base. Write
\[
 B=R^3,\qquad
 M=S^2=\{x=(x_1,x_2,x_3)\in R^3:x_1^2+x_2^2+x_3^2=1\},
 \qquad h(x)=x_1.
\]
Use base coordinates $(a,b,r)$ and define
\[
 G(a,b,r,x)=g(a,b,h(x)),\qquad
 Z=\{(a,b,r,x,w):x\in M,\ w=G(a,b,r,x)\}\subset R^7,
 \tag{MT.6}
\]
where $g$ is (MT.1). Let $p:Z\to B$ be the ambient linear
projection onto $(a,b,r)$. The estimate $|G|\le |b|/2$ proves
continuity. The graph map
\[
 B\times S^2\longrightarrow Z,\qquad
 (a,b,r,x)\longmapsto(a,b,r,x,G(a,b,r,x))
 \tag{MT.7}
\]
is a semialgebraic homeomorphism over $B$. The fibre is the closed
manifold $S^2$, the source is a topological $5$-manifold, and $p$
is definably proper: the graph over any closed bounded definable
$K\times S^2$ is closed and bounded. Its maximal open bundle
locus is all of $B$.

We prove that no neighbourhood of any base point $(0,0,r_0)$
admits a finite definable ambient $C^1$ Thom stratification for $p$.
Suppose such a stratification exists. Consider the central fibre
\[
 E_{r_0}=\{(0,0,r_0,x,0):x\in S^2\}.
\]
The open subset of this fibre where $|x_1|<1$ has dimension $2$.
Among its intersections with the finitely many source strata,
one has nonempty relative interior. Thus there is a source stratum
$T$, a point $x^0\in S^2$ with $|x^0_1|<1$, and a relative
neighbourhood of
$q=(0,0,r_0,x^0,0)$ in $E_{r_0}$ contained in $T$. Consequently
\[
 C=\{(0,0,0,v,0):v\in T_{x^0}S^2\}
       \subseteq\ker d(p|_T)_q.
 \tag{MT.8}
\]
Approach $q$ by
\[
 q_t=(t h(x^0),t,r_0,x^0,0),\qquad t>0.
 \tag{MT.9}
\]
This is a curve in $Z$. At every $q_t$, the source is a smooth
graph over an open subset of $B\times S^2$. Since
$g_z(t h(x^0),t,h(x^0))=1$, its full vertical tangent space is
the fixed $2$-plane
\[
 V=\{(0,0,0,v,(dh)_{x^0}v):v\in T_{x^0}S^2\}.
 \tag{MT.10}
\]
By o-minimality a tail of $q_t$ lies in one source stratum $S$.
The tangent kernels $K_t=\ker d(p|_S)_{q_t}$ are subspaces of
$V$ of constant dimension. They have a limit $L$ as $t\to0^+$:
one may represent them by their orthogonal projection matrices,
whose entries are bounded definable functions of one variable,
and take their limits. The limiting matrix is again a projection
of the same rank, so it represents a subspace $L\subseteq V$.
This argument applies over every real closed field; it uses
definable one-variable limits, not sequential compactness.

The vector
\[
 v_0=(1-(x^0_1)^2,-x^0_1x^0_2,-x^0_1x^0_3)
\]
belongs to $T_{x^0}S^2$, because its scalar product with $x^0$
is $x^0_1(1-\lVert x^0\rVert^2)=0$. It satisfies
$(dh)_{x^0}v_0=1-(x^0_1)^2>0$. Hence
$(0,0,0,v_0,0)\in C$ but $(0,0,0,v_0,0)\notin V$.
In particular, $C\not\subseteq V$.

If $S=T$, continuity of the constant-rank kernel bundle on $T$
would give $\ker d(p|_T)_q=L\subseteq V$, contradicting
(MT.8). Otherwise $q\in T\cap\overline S$, and the frontier
condition makes $S,T$ an incident pair. Thom's condition would
then give
\[
 C\subseteq\ker d(p|_T)_q\subseteq L\subseteq V,
\]
the same contradiction. This proves the claimed local
nonexistence for arbitrary finite compatible source and target
stratifications. In particular, restricting attention to closed
manifold fibres, a topological manifold source, target dimension
$3$, or fibre dimension $2$ does not repair the proposed
necessity of an ambient Thom stratification.

More generally, $S^2$ can be replaced by a closed definable
$C^1$ manifold $M$ carrying a definable $C^1$ function $h$ whose
differential is nonzero somewhere. Choose the central point in
the open regular locus of $h$; (MT.8)--(MT.10) give the same
argument with $T_{x^0}M$ in place of $T_{x^0}S^2$.

\paragraph{What this does and does not settle.}
This example disproves the proposed premise that a semialgebraically
locally trivial bundle is always locally Thom-stratifiable in its
given ambient smooth embedding, even after arbitrary finite
compatible refinements. It does not obstruct an intrinsic change
of differentiable structure: transporting the product structure by
$H$ makes $f$ a smooth product projection. It also does not refute
the original A.5 maximal-locus assertion: the locus here is $R^2$,
explicitly definable. Thus it is a failed-route certificate and
does not count as qualifying main-gate movement.

The separate elementary examples $x\mapsto x^2$ and $x\mapsto x^3$
distinguish two other failed criteria. Allowing target stratification,
$x^2$ is Thom-stratifiable but is not locally a product at $0$.
Requiring a submersion onto one full open target stratum excludes
the genuine product map $x^3$ at $0$. Neither convention alone
characterizes the desired topological/definable bundle locus.

The finite exact-algebra certificate is
\path{agents/maximal_locus/code/check_product_without_thom.py}.
Its initial interval-only run recorded ten passing checks in
$0.10073910000210162$ seconds. The extended version also checks
the sphere tangent-vector and height-differential identities;
the current exact output and runtime are recorded in
\path{agents/maximal_locus/logs/check_product_without_thom.json}.
It checks rational
identities, derivatives, the approach curve, and the two distinct
tangent lines, not an enumeration of possible stratifications.
We claim no priority. The construction was reached while checking
the product-bundle variant of the previous Thom-recognition route.
